% GENERATED FILE. Edit canonical lessons, then run npm run generate:books.
% canonical-source-hash: fnv1a64:f3b625c4ff564f0b
% canonical-lessons: SA-C296-davagnih, SA-C296-prajvalati, SA-C296-sphulingah, SA-C296-gadhah, SA-C296-prasarati

\chapter{Fire!}
\label{ch:sa-fire}
\hlchaptermodality{296}

\begin{chapteropening}
\textbf{By the end of this chapter:} \emph{I can say where a fire started and how it spread.}

Complete the last lesson of chapter 296: I can say where a fire started and how it spread.
\end{chapteropening}

\section[\texorpdfstring{\sk{दावाग्निः} (dāvāgniḥ)}{दावाग्निः (dāvāgniḥ)}]{\sk{दावाग्निः} (dāvāgniḥ) — a forest fire}

\label{lesson:SA-C296-davagnih}



\begin{quote}
Before the new one: say the Sanskrit for an emergency, then the Sanskrit for leans on.
\end{quote}

\subsection*{You'll want to know: \sk{दावाग्निः}}
\textbf{\sk{दावाग्निः}} — \emph{dāvāgniḥ} — \textquotedblleft{}a forest fire\textquotedblright{}.

\begin{quote}
\sk{पर्वते} \sk{दावाग्निः} \sk{ज्वलति।}
\end{quote}

\begin{quote}
\emph{parvate dāvāgniḥ jvalati.}
\end{quote}

\begin{quote}
A forest fire is burning on the hill.
\end{quote}

A masculine noun.

\begin{cousinweb}
\textbf{\sk{दावाग्निः}} (\emph{dāvāgniḥ}) is built from \emph{dāva}, \textquotedblleft{}a blaze, a forest blaze\textquotedblright{}, and \emph{agni}, \textquotedblleft{}fire\textquotedblright{}; MW gives \emph{dāvāgni} as fire in a forest (MW).
\end{cousinweb}

\begin{grammarlens}[title={What you've built}]
One more word for telling what happened.
\end{grammarlens}

\subsection*{Your turn}
\begin{itemize}
\raggedright
  \item \emph{Say it:} \emph{dāvāgniḥ}
  \item \emph{Say it:} \emph{dāvāgniḥ}, once more
  \item \emph{Say it:} \emph{parvate dāvāgniḥ jvalati.} — A forest fire is burning on the hill.
  \item \emph{Recall:} say the Sanskrit for an emergency, then the Sanskrit for leans on, then say \emph{dāvāgniḥ} again
\end{itemize}

\begin{tcolorbox}[breakable,colback=teal!4,colframe=teal!35!black,title={Before you move on}]
What does \sk{दावाग्निः} mean? (\textquotedblleft{}A forest fire\textquotedblright{}.) Say it once more.
\end{tcolorbox}
\section[\texorpdfstring{\sk{प्रज्वलति} (prajvalati)}{प्रज्वलति (prajvalati)}]{\sk{प्रज्वलति} (prajvalati) — blazes up; catches fire}

\label{lesson:SA-C296-prajvalati}



\begin{quote}
Before the new one: say the Sanskrit for leans on, then the Sanskrit for a forest fire.
\end{quote}

\subsection*{You'll want to know: \sk{प्रज्वलति}}
\textbf{\sk{प्रज्वलति}} — \emph{prajvalati} — \textquotedblleft{}blazes up; catches fire\textquotedblright{}.

\begin{quote}
\sk{शुष्कं} \sk{तृणं} \sk{सहसा} \sk{प्रज्वलति।}
\end{quote}

\begin{quote}
\emph{śuṣkaṃ tṛṇaṃ sahasā prajvalati.}
\end{quote}

\begin{quote}
The dry grass suddenly blazes up.
\end{quote}

Past (it …): \sk{तृणं} \sk{सहसा} \sk{प्रज्वलितवत्}, \emph{tṛṇaṃ sahasā prajvalitavat}.

\begin{cousinweb}
\textbf{\sk{प्रज्वलति}} (\emph{prajvalati}) is built from \emph{pra}, \textquotedblleft{}forth\textquotedblright{}, and the root √jval, \textquotedblleft{}to burn, blaze\textquotedblright{}; MW gives \emph{pra-jval} as to begin to burn, to flame up (MW).
\end{cousinweb}

\begin{grammarlens}[title={What you've built}]
One more word for telling what happened.
\end{grammarlens}

\subsection*{Your turn}
\begin{itemize}
\raggedright
  \item \emph{Say it:} \emph{prajvalati}
  \item \emph{Say it:} \emph{prajvalati}, once more
  \item \emph{Say it:} \emph{śuṣkaṃ tṛṇaṃ sahasā prajvalati.} — The dry grass suddenly blazes up.
  \item \emph{Recall:} say the Sanskrit for leans on, then the Sanskrit for a forest fire, then say \emph{prajvalati} again
\end{itemize}

\begin{tcolorbox}[breakable,colback=teal!4,colframe=teal!35!black,title={Before you move on}]
What does \sk{प्रज्वलति} mean? (\textquotedblleft{}Blazes up; catches fire\textquotedblright{}.) Say it once more.
\end{tcolorbox}
\section[\texorpdfstring{\sk{स्फुलिंगः} (sphuliṅgaḥ)}{स्फुलिंगः (sphuliṅgaḥ)}]{\sk{स्फुलिंगः} (sphuliṅgaḥ) — a spark}

\label{lesson:SA-C296-sphulingah}



\begin{quote}
Before the new one: say the Sanskrit for a forest fire, then the Sanskrit for blazes up.
\end{quote}

\subsection*{You'll want to know: \sk{स्फुलिंगः}}
\textbf{\sk{स्फुलिंगः}} — \emph{sphuliṅgaḥ} — \textquotedblleft{}a spark\textquotedblright{}.

\begin{quote}
\sk{अग्नेः} \sk{स्फुलिंगाः} \sk{वायुना} \sk{उड्डयन्ते।}
\end{quote}

\begin{quote}
\emph{agneḥ sphuliṅgāḥ vāyunā uḍḍayante.}
\end{quote}

\begin{quote}
Sparks from the fire fly on the wind.
\end{quote}

A masculine noun.

\begin{cousinweb}
\textbf{\sk{स्फुलिंगः}} (\emph{sphuliṅgaḥ}) MW gives \emph{sphuliṅga} as a spark of fire and suggests, with a perhaps, that it is built from \emph{sphulam} and \emph{ga}, so its origin is uncertain (MW).
\end{cousinweb}

\begin{grammarlens}[title={What you've built}]
One more word for telling what happened.
\end{grammarlens}

\subsection*{Your turn}
\begin{itemize}
\raggedright
  \item \emph{Say it:} \emph{sphuliṅgaḥ}
  \item \emph{Say it:} \emph{sphuliṅgaḥ}, once more
  \item \emph{Say it:} \emph{agneḥ sphuliṅgāḥ vāyunā uḍḍayante.} — Sparks from the fire fly on the wind.
  \item \emph{Recall:} say the Sanskrit for a forest fire, then the Sanskrit for blazes up, then say \emph{sphuliṅgaḥ} again
\end{itemize}

\begin{tcolorbox}[breakable,colback=teal!4,colframe=teal!35!black,title={Before you move on}]
What does \sk{स्फुलिंगः} mean? (\textquotedblleft{}A spark\textquotedblright{}.) Say it once more.
\end{tcolorbox}
\section[\texorpdfstring{\sk{गाढः} (gāḍhaḥ)}{गाढः (gāḍhaḥ)}]{\sk{गाढः} (gāḍhaḥ) — dense; thick}

\label{lesson:SA-C296-gadhah}



\begin{quote}
Before the new one: say the Sanskrit for blazes up, then the Sanskrit for a spark.
\end{quote}

\subsection*{You'll want to know: \sk{गाढः}}
\textbf{\sk{गाढः}} — \emph{gāḍhaḥ} — \textquotedblleft{}dense; thick\textquotedblright{}.

\begin{quote}
\sk{सर्वत्र} \sk{गाढः} \sk{धूमः} \sk{अस्ति।}
\end{quote}

\begin{quote}
\emph{sarvatra gāḍhaḥ dhūmaḥ asti.}
\end{quote}

\begin{quote}
There is thick smoke everywhere.
\end{quote}

\begin{cousinweb}
\textbf{\sk{गाढः}} (\emph{gāḍhaḥ}) comes from the root √gāh, \textquotedblleft{}to plunge into\textquotedblright{}: MW gives \emph{gāḍha} as deeply entered, close, and also thick or dense (MW).
\end{cousinweb}

\begin{grammarlens}[title={What you've built}]
One more word for telling what happened.
\end{grammarlens}

\subsection*{Your turn}
\begin{itemize}
\raggedright
  \item \emph{Say it:} \emph{gāḍhaḥ}
  \item \emph{Say it:} \emph{gāḍhaḥ}, once more
  \item \emph{Say it:} \emph{sarvatra gāḍhaḥ dhūmaḥ asti.} — There is thick smoke everywhere.
  \item \emph{Recall:} say the Sanskrit for blazes up, then the Sanskrit for a spark, then say \emph{gāḍhaḥ} again
\end{itemize}

\begin{tcolorbox}[breakable,colback=teal!4,colframe=teal!35!black,title={Before you move on}]
What does \sk{गाढः} mean? (\textquotedblleft{}Dense; thick\textquotedblright{}.) Say it once more.
\end{tcolorbox}
\section[\texorpdfstring{\sk{प्रसरति} (prasarati)}{प्रसरति (prasarati)}]{\sk{प्रसरति} (prasarati) — spreads}

\label{lesson:SA-C296-prasarati}



\begin{quote}
Before the new one: say the Sanskrit for a spark, then the Sanskrit for dense.
\end{quote}

\subsection*{You'll want to know: \sk{प्रसरति}}
\textbf{\sk{प्रसरति}} — \emph{prasarati} — \textquotedblleft{}spreads\textquotedblright{}.

\begin{quote}
\sk{वायुना} \sk{अग्निः} \sk{शीघ्रं} \sk{प्रसरति।}
\end{quote}

\begin{quote}
\emph{vāyunā agniḥ śīghraṃ prasarati.}
\end{quote}

\begin{quote}
With the wind the fire spreads quickly.
\end{quote}

Past (it …): \sk{अग्निः} \sk{सर्वत्र} \sk{प्रसृतवान्}, \emph{agniḥ sarvatra prasṛtavān}.

\begin{cousinweb}
\textbf{\sk{प्रसरति}} (\emph{prasarati}) is built from \emph{pra}, \textquotedblleft{}forth\textquotedblright{}, and the root √sṛ, \textquotedblleft{}to run, flow\textquotedblright{}; MW gives \emph{pra-sṛ} as to move forwards and spread, and to break out, as a fire does (MW).
\end{cousinweb}

\begin{grammarlens}[title={What you've built}]
One more word for telling what happened.
\end{grammarlens}

\subsection*{Your turn}
\begin{itemize}
\raggedright
  \item \emph{Say it:} \emph{prasarati}
  \item \emph{Say it:} \emph{prasarati}, once more
  \item \emph{Say it:} \emph{vāyunā agniḥ śīghraṃ prasarati.} — With the wind the fire spreads quickly.
  \item \emph{Recall:} say the Sanskrit for a spark, then the Sanskrit for dense, then say \emph{prasarati} again
\end{itemize}

\begin{tcolorbox}[breakable,colback=teal!4,colframe=teal!35!black,title={Before you move on}]
What does \sk{प्रसरति} mean? (\textquotedblleft{}Spreads\textquotedblright{}.) Say it once more.
\end{tcolorbox}
